\chapter{A new robot}
\label{ch:i-bringup}

Before the stack walks a robot on its own, check that the robot does what the stack tells it,
and measure what the stack assumes about it. In an open area, one person at the handheld whose
only job is to stop the robot.

\section{Does the robot obey?}

Stop the service (\cmd{sudo systemctl stop x30-navigation}) and start only the driver:

\begin{lstlisting}[language=bash]
eval "$(bash ~/x30_ws/src/x30_ros/tools/x30_host.sh env navigation)"
roslaunch x30_bringup bringup.launch
\end{lstlisting}

\begin{center}
\begin{small}
\begin{longtable}{@{}p{0.04\linewidth}>{\raggedright}p{0.38\linewidth}>{\raggedright\arraybackslash}p{0.5\linewidth}@{}}
  \toprule
  & \textbf{Do} & \textbf{Expect} \\
  \midrule
  \endhead
  \bottomrule
  \endfoot
  1 & \cmd{rostopic hz /x30/odom} & about 50~Hz; \topic{/x30/state} says \code{connected: True}.
    If not: \file{network.toml} (Section~\ref{sec:i-networktoml}) \\
  2 & \cmd{rostopic echo -n 1 /x30/state} & \code{rcs_layout} and \code{motion_layout} read
    \code{natural} \\
  3 & \cmd{rosservice call /x30/stand_up} & ends in \code{torque_standing};
    \cmd{/x30/start_stepping} in \code{stepping} \\
  4 & \cmd{rostopic pub -r 20 /x30/cmd_vel geometry_msgs/Twist '{linear: {x: 0.2}}'} & the
    robot walks forward, \code{control_mode} turns \code{non_manual}; stop the publisher: it
    stops within 0.5~s \\
  5 & the same with \code{{angular: {z: 0.3}}} & it turns \textbf{left}
    (counter-clockwise seen from above); else \code{invert_wz: true} \\
  6 & the same with \code{{linear: {y: 0.1}}} & it steps \textbf{left}; else
    \code{invert_vy: true} \\
  7 & walking, switch to manual on Deep Robotics' handheld & the robot follows that handheld;
    \code{velocity} reads \code{takeover (manual mode)} until it is handed back \\
  8 & \cmd{rostopic echo -n 1 /x30/battery} & voltage and current as Deep Robotics' app shows
    them \\
  9 & \cmd{roslaunch x30_bringup bringup.launch rviz:=true} & the legs in RViz follow the real ones \\
\end{longtable}
\end{small}
\end{center}

A robot that disagrees with step 5 or 6 gets its own copy of
\file{x30_driver/config/driver_params.yaml} with the inversion, passed as
\code{params_file:=}. The Tuning Guide explains every driver setting.

\section{The robot profile}
\label{sec:i-profile}

\file{x30_bringup/config/robot.yaml} is the X30 for every node: its footprint, its height, its
limits and its frames. Three things are particular to each robot and are measured, not taken
from the drawing:

\begin{enumerate}
  \item \textbf{\code{base_height}}: \code{base_link} (the body centre) above the floor while it
    stands. The scans measure their height band from the floor \code{base_height} below it;
    0.48~m is the model's.
  \item \textbf{The LiDAR mount}: \cmd{rostopic echo -n 1 --noarr /livox/lidar/header} gives the
    frame the robot's driver publishes in. Set \code{lidar_frame} and \code{lidar_pose} so that a
    wall 1~m ahead shows at 1~m in \cmd{rostopic echo -n 1 /perception/safety_scan}.
  \item \textbf{Braking}: \code{limits/max_decel}, \code{max_yaw_accel} and
    \code{reaction_time}, worst case (floor, payload, low battery). The protective field is sized
    from them; the Tuning Guide shows how to measure them.
\end{enumerate}

\section{The first site and route}

With the services running again (\cmd{sudo systemctl start x30-navigation}):

\begin{enumerate}
  \item Put the robot on its start spot (the dock), facing the same way every time: the start of
    the odometry is the site's origin, and one degree off there is 1.7~cm per metre on every route.
  \item \cmd{rostopic echo -n 1 /localization/status}: \code{LOCALIZED}, source \code{origin}.
  \item Teach a short route with the H12 (User Manual): REMOTE, \ui{Record}, walk, \ui{Point} at
    a stop, \ui{Finish}.
  \item Walk back to the start, pick the route in \ui{Routes}, AUTO, \ui{Start}. The robot repeats
    it and stops at the stop point.
  \item In NETRA the route appears under \menupath{Routes}, and the robot's patrol state on the
    dashboard.
\end{enumerate}

Then the acceptance tests (X30 Site Acceptance Test) before the robot patrols for real.
